\documentclass[UTF8,11pt,a4paper]{ctexart}
\usepackage[margin=24mm]{geometry}
\usepackage{amsmath,amssymb,amsthm,mathtools,bm}
\usepackage{tikz-cd,enumitem,longtable,booktabs}
\usepackage{xurl}
\usepackage[unicode,hidelinks]{hyperref}
\allowdisplaybreaks
\setlength{\emergencystretch}{3em}
\setlength{\parindent}{0pt}
\setlength{\parskip}{0.45em}
\newtheorem*{theorem}{Theorem}
\newtheorem*{lemma}{Lemma}
\newtheorem*{proposition}{Proposition}
\newtheorem*{corollary}{Corollary}
\theoremstyle{definition}\newtheorem*{definition}{Definition}
\theoremstyle{remark}\newtheorem*{remark}{Remark}
\DeclareMathOperator{\Hom}{Hom}
\DeclareMathOperator{\End}{End}
\DeclareMathOperator{\Tr}{Tr}
\DeclareMathOperator{\rank}{rank}
\DeclareMathOperator{\Ind}{Ind}
\DeclareMathOperator{\Res}{Res}
\DeclareMathOperator{\im}{Im}
\DeclareMathOperator{\id}{id}
\newcommand{\R}{\mathbb R}
\newcommand{\C}{\mathbb C}
\newcommand{\Q}{\mathbb Q}
\newcommand{\Z}{\mathbb Z}
\newcommand{\N}{\mathbb N}
\newcommand{\E}{\mathbb E}
\newcommand{\Prob}{\mathbb P}
\newcommand{\eps}{\varepsilon}
\begin{document}
\section*{098｜仅解码端有侧信息时的Slepian--Wolf压缩极限}
\textbf{作者／集体来源：}Yury Polyanskiy and Yihong Wu

\textbf{作品：}MIT 6.441 Information Theory, Spring 2016, Chapter 7

\textbf{原文入口：}\url{https://ocw.mit.edu/courses/6-441-information-theory-spring-2016/85014cb1add9f415ef38f60f678f16b2_MIT6_441S16_chapter_7.pdf}

\textbf{原文位置与计题范围：}§7.4，Definition 7.4、Theorem 7.7、Corollary 7.3（页84--85）；附§7.3条件化论证及页77 Theorem 7.2的逆定理。单计侧信息只在解码端可见的压缩问题。

\textbf{获取日期：}2026-09-28。\quad\textbf{复用依据：}MIT OCW CC BY-NC-SA 4.0。

\textbf{整理归属：}下列题面与证明取自所列人类来源，按注明范围转录、摘编；LaTeX环境、标题、换行及中文说明由助手整理。编辑调整在题内说明。

\section*{作者命题与人类证明}

\paragraph{Notation and scope (editorial).}
Logarithms in this item are base $2$. The source $X$ is discrete; $Y$ is side information. $\eps^*(X\mid Y,k)$ denotes the optimal error probability when both encoder and decoder observe $Y$; $\eps^*_{\rm SW}(X\mid Y,k)$ denotes the optimal error probability when only the decoder observes $Y$. The i.i.d. statement assumes finite conditional entropy. An integer block length $k$ is understood in expressions such as $k=nR$ (rounding does not change the limits).

\subsection*{Converse bound retained from Theorem 7.2}
\begin{theorem}[Converse, 7.2]
For compression of $X$ into $k$ bits, with undetectable errors allowed,
\[
\widetilde\eps^*(X,k)\ge
\Prob\left[\log\frac1{P_X(X)}>k+\tau\right]-2^{-\tau},\qquad\tau>0.
\]
\end{theorem}
\begin{proof}
Let $C=\{x:g(f(x))=x\}$. Then $|C|\le2^k$ and
\[
\Prob[X\in C]\le
\Prob\left[\log\frac1{P_X(X)}\le k+\tau\right]
+\underbrace{\Prob\left[X\in C,\log\frac1{P_X(X)}>k+\tau\right]}_{\le2^{-\tau}}.
\]
\end{proof}

\subsection*{Conditioning the bound: source \S7.3}
Conditioned on $Y=y$, the problem reduces to compression without side information, where the source $X$ is distributed according to $P_{X\mid Y=y}$. Since $Y$ is known to both the compressor and decompressor, they can use the best code tailored for this distribution.
\[
\eps^*(X\mid Y,k)=\E_{y\sim P_Y}[\eps^*(P_{X\mid Y=y},k)].
\]
Using the converse Theorem 7.2 for compression without side information and taking the average over all $y\sim P_Y$ gives
\[
\Prob\left[\log\frac1{P_{X\mid Y}(X\mid Y)}>k+\tau\right]-2^{-\tau}
\le\eps^*(X\mid Y,k).
\]
For $(X,Y)=(S^n,T^n)$ where $(S_1,T_1),(S_2,T_2),\ldots$ are iid pairs $\sim P_{ST}$,
\[
\frac1n\log\frac1{P_{S^n\mid T^n}(S^n\mid T^n)}
=\frac1n\sum_{i=1}^n\log\frac1{P_{S\mid T}(S_i\mid T_i)}
\longrightarrow H(S\mid T)\quad\text{in probability}
\]
as $n\to\infty$, using the WLLN.

\subsection*{Slepian--Wolf: compression with side information at decompressor only}
Consider the compression with side information problem, except now the compressor has no access to the side information.
\begin{center}
\begin{tikzcd}[column sep=large,row sep=small]
X\arrow[r]&f(X)\in\{0,1\}^k\arrow[r]&g(f(X),Y)\in\mathcal X\cup\{e\}\\
&&Y\arrow[u]
\end{tikzcd}
\end{center}
\begin{definition}[S.W. code, 7.4]
Given $P_{XY}$, a $(k,\eps)$-S.W. code consists of
\[
f:\mathcal X\to\{0,1\}^k,\qquad
g:\{0,1\}^k\times\mathcal Y\to\mathcal X\cup\{e\},
\]
with $\Prob[g(f(X),Y)\ne X]\le\eps$. The fundamental limit is
\[
\eps^*_{\rm SW}(X\mid Y,k)=\inf\{\eps:\exists\,(k,\eps)\text{-S.W. code}\}.
\]
\end{definition}
\begin{theorem}[Slepian--Wolf, 1973; 7.7]
\[
\eps^*(X\mid Y,k)\le\eps^*_{\rm SW}(X\mid Y,k)
\le\Prob\left[\log\frac1{P_{X\mid Y}(X\mid Y)}\ge k-\tau\right]+2^{-\tau}.
\]
\end{theorem}
\begin{corollary}[7.3]
\[
\lim_{n\to\infty}\eps^*_{\rm SW}(S^n\mid T^n,nR)
=\begin{cases}0,&R>H(S\mid T),\\1,&R<H(S\mid T).
\end{cases}
\]
\end{corollary}
\paragraph{Source's qualification.}
Definition 7.4 does not include the zero-undetected-error condition (that is $g(f(x),y)=x$ or $e$). In other words, we allow for the possibility of undetected errors.
\begin{proof}[Proof of Theorem 7.7]
LHS is obvious, since side information at the compressor and decoder is better than only at the decoder.

For the RHS, first generate a random codebook with iid uniform codewords:
$C=\{c_x\in\{0,1\}^k:x\in\mathcal X\}$ independently of $(X,Y)$, then define the compressor and decoder as
\[
f(x)=C_x,\qquad
 g(w,y)=\begin{cases}x,&\exists!x:C_x=w,\ x\text{--h.p.}\mid y,\\0,&\text{o.w.}\end{cases}
\]
where we used the shorthand
\[
x\text{--h.p.}\mid y\quad\Longleftrightarrow\quad
\log\frac1{P_{X\mid Y}(x\mid y)}<k-\tau.
\]
The error probability of this scheme is bounded by
\begin{align*}
\mathcal E(C)
&=\Prob\left[\log\frac1{P_{X\mid Y}(X\mid Y)}\ge k-\tau
\ \text{or }J(X,C\mid Y)\ne\varnothing\right]\\
&\le\Prob\left[\log\frac1{P_{X\mid Y}(X\mid Y)}\ge k-\tau\right]
+\Prob[J(X,C\mid Y)\ne\varnothing]\\
&=\Prob\left[\log\frac1{P_{X\mid Y}(X\mid Y)}\ge k-\tau\right]
+\sum_{x,y}P_{XY}(x,y)\mathbf1_{\{J(x,C\mid y)\ne\varnothing\}},
\end{align*}
where $J(x,C\mid y)\mathrel{:=}\{x'\ne x:x'\text{--h.p.}\mid y,\ c_x=c_{x'}\}$.
Now averaging over $C$ and applying the union bound: use
$|\{x':x'\text{--h.p.}\mid y\}|\le2^{k-\tau}$ and
$\Prob[C_{x'}=C_x]=2^{-k}$ for any $x\ne x'$,
\begin{align*}
\Prob_C[J(x,C\mid y)\ne\varnothing]
&\le\E_C\left[\sum_{x'\ne x}
\mathbf1_{\{x'\text{--h.p.}\mid y\}}\mathbf1_{\{C_{x'}=C_x\}}\right]\\
&\le2^{k-\tau}\Prob[C_{x'}=C_x]=2^{-\tau}.
\end{align*}
Hence the theorem follows as usual from two terms in the union bound.
\end{proof}
\paragraph{Connection to Corollary 7.3.}
The i.i.d. reduction is the conditional-information WLLN displayed in the preceding source proof (Corollary 7.2). Combined with the two bounds above, it gives the two sides of Corollary 7.3. No claim is made at $R=H(S\mid T)$.

\section*{整理说明与可修改标注（助手）}
\textbf{标准先修：}离散概率、条件熵、联合上界、弱大数定律；辅助的有限码本逆界已附入。

\textbf{本题仍需完成的工作：}构造与源独立的随机码本，仅在解码时借助条件高概率集消除碰撞，并结合条件信息密度得到可达性与强逆结论。

\textbf{取材说明：}记号与流程图由助手按原文整理。原文页85将坏事件概率写成实际误码率的等号，本稿明确记为误码率上界；计数估计中原印等号改为小于等于。解码器的默认0沿用作者记号。§7.3只摘取所需的条件化逆界与WLLN段，未声称另录完整的双端侧信息定理；最后衔接段由助手说明原文7.3与此前WLLN的对应。原页77、83--85已查看。

\textbf{$c$：}侧信息约束下的近无损压缩

\textbf{$a$：}随机编码；条件高概率集；碰撞概率；条件信息密度；大数定律

\texttt{cross\_theory\_status = yes}

\textbf{具体联系及位置：}页85将解码失败拆成条件信息密度尾部和随机码本碰撞两部分；码本不依赖Y而候选集依赖Y，正是两端信息不对称下达到条件熵的关键。

这些标注是非穷尽初稿；未列出不表示未使用。
\end{document}
